% answers_short.tex —— 2010 年·普通高等学校招生考试·福建卷·理科 · 简短答案速查
% 由 tools/make_short_answers.py 自 answers.tex 的 \gksolhead/\gkanswer 提取，
% 并与 2010/2010_数学真题_福建理_答案要点.md 逐题比对后生成，勿手改（改源 MD 后重跑）。

\begin{gknotice}
\textbf{使用说明}：本节只给最终答案便于快速核对，完整推导见解析版。答案为编写者独立推导，\textbf{非官方标准答案}。
\end{gknotice}

\gkpart{选择题}
\begin{center}\begin{tabular}{ccccccccccc}
\toprule
题号 & 1 & 2 & 3 & 4 & 5 & 6 & 7 & 8 & 9 & 10 \\
\midrule
答案 & A & D & A & C & C & D & B & B & B & C \\
\bottomrule
\end{tabular}\end{center}

\gkpart{填空题}
\begin{itemize}[leftmargin=2.6em]
  \item[\textbf{11.}] $a_n=4^{n-1}$。
  \item[\textbf{12.}] $6+2\sqrt3$。
  \item[\textbf{13.}] $0.128$。
  \item[\textbf{14.}] $\left[-\dfrac32,3\right]$。
  \item[\textbf{15.}] ①②④。
\end{itemize}

\gkpart{解答题}
\begin{itemize}[leftmargin=2.6em]
  \item[\textbf{16.}] （1）$(-2,2)$，$(-1,1)$，$(0,0)$，$(1,-1)$，$(2,-2)$，共 $5$ 个基本事件；（2）$P(\xi=0)=\dfrac16$，$P(\xi=1)=\dfrac13$，$P(\xi=4)=\dfrac13$，$P(\xi=9)=\dfrac16$，$E\xi=\dfrac{19}{6}$。
  \item[\textbf{17.}] （1）$\dfrac{x^2}{16}+\dfrac{y^2}{12}=1$；（2）不存在。
  \item[\textbf{18.}] （1）证明见解析；（2）① $p_{\max}=\dfrac1\pi$；② $\cos\theta=\dfrac{\sqrt{10}}{5}$。
  \item[\textbf{19.}] （1）$30\sqrt3$ 海里/小时；（2）沿北偏东 $30^\circ$ 方向、以 $30$ 海里/小时航行，最短时间 $\dfrac23$ 小时。
  \item[\textbf{20.}] （1）① 增区间 $\left(-\infty,-\dfrac{\sqrt3}{3}\right)$，$\left(\dfrac{\sqrt3}{3},+\infty\right)$，减区间 $\left(-\dfrac{\sqrt3}{3},\dfrac{\sqrt3}{3}\right)$；② $\dfrac{S_1}{S_2}=\dfrac1{16}$ 为定值；（2）命题见解析，$\dfrac{S_1}{S_2}=\dfrac1{16}$ 仍为定值。
  \item[\textbf{21.}] 【A】$a=b=-1$，$c=d=2$；像的方程为 $y=-x$。【B】$x^2+(y-\sqrt5)^2=5$；$|PA|+|PB|=3\sqrt2$。【C】$a=2$；$m\leqslant5$。
\end{itemize}
